\uuid{KZPz}
\titre{Domaine de définition}
\auteur{bodin ; hanani ; mzari}
\organisation{exo7}
\licence{CC BY-NC-SA 4.0}
\datecreate{2020-01-20}
\niveau{L1}
\module{Analyse}
\chapitre{Fonctions circulaires et hyperboliques inverses}
\sousChapitre{Autre}
\difficulte{1}
\theme{Fonctions usuelles, domaine de définition}
\exo7qcmid{lille-1.fonctions-usuelles.007}

\contenu{
\begin{question}
\qclassification{126.00}
\qtags{motcle=domaine de définition}
Soit $(E)$ l'équation : $\ln (x^2-1) = \ln (x-1) + \ln 2$. Quelles sont les assertions vraies ?
\begin{answers}
    \bad{$(E)$ est définie sur  $]-\infty, -1[ \cup  ]1,+\infty[$.}

    \good{$(E)$ est définie sur  $]1,+\infty[$.}

    \good{$(E)$ n'admet pas de solution.}

    \bad{$(E)$ admet une unique solution $x=1$.}
\end{answers}
\begin{explanations}
$(E)$ est définie si $x^2-1>0$ et $x-1>0$, c'est-à-dire $x>1$.
Soit $x>1$, alors $(E) \Leftrightarrow  \ln (x-1) + \ln (x+1)= \ln (x-1) + \ln 2  \Leftrightarrow  \ln (x+1)=  \ln 2 \Leftrightarrow  x=1$.
\end{explanations}
\end{question}
}
